%%%%%%%%% ELSEVIER JOURNAL TEMPLATE (elsarticle class) %%%%%%%%%
% Journal: Computational Materials Science (Elsevier). For another Elsevier
% journal, change only the \journal{...} line below (see guidelines.md).
% Class: elsarticle v3.3 (elsarticle.cls and elsarticle-num.bst are copied
% next to this file, see SOURCES.md).
%
% Submission manuscript: single column, 12 pt, A4 (Elsevier's default).
\documentclass[preprint,12pt]{elsarticle}
% Two-column look of the final journal pages (preview only, not for
% submission): comment out the line above and uncomment this one.
% \documentclass[final,5p,times,twocolumn]{elsarticle}

%%%%%%%%% PACKAGES - KEEP UNCHANGED %%%%%%%%%
% elsarticle.cls already loads graphicx and natbib (numbered citations in
% square brackets) and sets the page layout. Do not load natbib, cite,
% geometry, hyperref, cleveref or a font package.
\usepackage{amsmath}
\usepackage{amssymb}

% Tables: booktabs (\toprule, \midrule, \cmidrule, \bottomrule).
\IfFileExists{booktabs.sty}{\usepackage{booktabs}}{}
\makeatletter
\providecommand*{\toprule}{\hline}
\providecommand*{\midrule}{\hline}
\providecommand*{\bottomrule}{\hline}
\providecommand*{\addlinespace}[1][]{}
\providecommand*{\cmidrule}{\@ifnextchar(\els@po@cmidtrim\els@po@cmid}
\def\els@po@cmidtrim(#1){\els@po@cmid}
\def\els@po@cmid#1{\cline{#1}}
\makeatother

% URLs in the Data availability statement: \url{https://...}.
\IfFileExists{url.sty}{\usepackage{url}}{}
\providecommand*{\url}[1]{\texttt{\detokenize{#1}}}

% Chemistry: mhchem (version 4) for formulas and reactions, \ce{...}.
\IfFileExists{mhchem.sty}{\usepackage[version=4]{mhchem}}{}
% Fallback when mhchem is missing: print the formula as typed.
\providecommand*{\ce}[1]{\texttt{\detokenize{#1}}}

% Quantities and units: siunitx. \qty{1.5}{\kilo\calorie\per\mole} (or \SI),
% \unit{\electronvolt\per\atom} (or \si), \num{1.2e-3}.
\IfFileExists{siunitx.sty}{\usepackage{siunitx}}{}
\makeatletter
% siunitx has no \calorie and no \atom; declare them (and \angstrom).
\@ifpackageloaded{siunitx}{%
  \DeclareSIUnit{\calorie}{cal}%
  \DeclareSIUnit{\angstrom}{\text{\AA}}%
  \DeclareSIUnit{\atom}{atom}%
}{}
% Fallbacks when siunitx is missing (or is version 2, without \qty/\unit):
% common prefixes and units are printed upright.
\newcommand*{\els@po@units}{%
  \def\per{/}\def\squared{^{2}}\def\cubed{^{3}}\def\tothe##1{^{##1}}%
  \def\square##1{##1^{2}}\def\cubic##1{##1^{3}}%
  \def\femto{f}\def\pico{p}\def\nano{n}\def\micro{\mu}\def\milli{m}%
  \def\centi{c}\def\deci{d}\def\kilo{k}\def\mega{M}\def\giga{G}%
  \def\tera{T}\def\metre{m}\def\meter{m}\def\gram{g}\def\second{s}%
  \def\mole{mol}\def\kelvin{K}\def\joule{J}\def\calorie{cal}%
  \def\electronvolt{eV}\def\hartree{E_{\mathrm{h}}}\def\bohr{a_{0}}%
  \def\angstrom{\mbox{\AA}}\def\atom{atom}\def\dalton{Da}%
  \def\atomicmassunit{u}\def\litre{L}\def\liter{L}\def\hour{h}%
  \def\minute{min}\def\pascal{Pa}\def\bar{bar}\def\atmosphere{atm}%
  \def\hertz{Hz}\def\watt{W}\def\volt{V}\def\ampere{A}\def\coulomb{C}%
  \def\newton{N}\def\ohm{\Omega}\def\siemens{S}\def\tesla{T}%
  \def\debye{D}\def\molar{M}\def\percent{\%}\def\degree{^{\circ}}%
  \def\degreeCelsius{{}^{\circ}C}%
  \def\eV{eV}\def\meV{meV}\def\kJ{kJ}\def\J{J}\def\mol{mol}\def\K{K}%
  \def\s{s}\def\ns{ns}\def\ps{ps}\def\fs{fs}\def\nm{nm}\def\cm{cm}%
  \def\mm{mm}\def\kg{kg}\def\g{g}\def\Pa{Pa}\def\kPa{kPa}\def\MPa{MPa}%
  \def\GPa{GPa}\def\Hz{Hz}\def\THz{THz}%
}
\providecommand*{\si}[2][]{{\els@po@units\ensuremath{\mathrm{#2}}}}
\makeatother
\providecommand*{\num}[2][]{#2}
\providecommand*{\SI}[3][]{\num{#2}\,\si{#3}}
\providecommand*{\SIrange}[4][]{\num{#2}--\num{#3}\,\si{#4}}
\providecommand*{\numrange}[3][]{\num{#2}--\num{#3}}
\providecommand*{\qty}{\SI}
\providecommand*{\unit}{\si}
\providecommand*{\qtyrange}{\SIrange}

%%%%%%%%% JOURNAL - change for another Elsevier journal %%%%%%%%%
\journal{Computational Materials Science}

\begin{document}

\begin{frontmatter}

%%%%%%%%% TITLE - PLEASE UPDATE
\title{My Title}

%%%%%%%%% AUTHOR - KEEP THIS PLACEHOLDER EXACTLY AS IT IS
% elsarticle macros: \corref and \cortext mark the corresponding author
% (asterisk footnote), \ead prints the e-mail address in the footnote,
% and the label [inst1] links the author to the affiliation (superscript a).
\author[inst1]{A. Researcher\corref{cor1}}
\ead{researcher@example.edu}
\cortext[cor1]{Corresponding author.}
\affiliation[inst1]{organization={Department of Chemistry, Example University},
            city={City},
            country={Country}}

%%%%%%%%% GRAPHICAL ABSTRACT %%%%%%%%%
% Required by the journal. The class prints it on its own page before the
% title page. Elsevier size: readable at 13 cm x 5 cm (w x h), image at least
% 1328 x 531 pixels (w x h). To fill it, replace both occurrences of
% figures/toc_graphic.png with the most representative figure in figures/
% (see guidelines.md); keep the placeholder branch.
\begin{graphicalabstract}
\IfFileExists{figures/toc_graphic.png}{%
  {\centering
  \includegraphics[width=13cm,height=5cm,keepaspectratio]{figures/toc_graphic.png}\par}%
}{%
  {\centering
  \fbox{\parbox[c][5cm][c]{12.5cm}{\centering
    Graphical abstract placeholder\\[4pt]
    13~cm $\times$ 5~cm (w $\times$ h), at least 1328 $\times$ 531 pixels}}\par}%
}
\end{graphicalabstract}
%%%%%%%%% GRAPHICAL ABSTRACT %%%%%%%%%

%%%%%%%%% HIGHLIGHTS %%%%%%%%%
% Required by the journal. The class prints them on their own page before
% the title page. 3 to 5 \item lines, each at most 85 characters including
% spaces; no citations, no undefined abbreviations.
\begin{highlights}
\item Highlights here.
\end{highlights}
%%%%%%%%% HIGHLIGHTS %%%%%%%%%

\begin{abstract}
%%%%%%%%% ABSTRACT %%%%%%%%%
Abstract here.
%%%%%%%%% ABSTRACT %%%%%%%%%
\end{abstract}

\begin{keyword}
%%%%%%%%% KEYWORDS %%%%%%%%%
% 1 to 7 keywords, separated by \sep (e.g. Density functional theory \sep
% Phonons \sep Thermal conductivity).
Keywords here.
%%%%%%%%% KEYWORDS %%%%%%%%%
\end{keyword}

\end{frontmatter}
% The class prints the title block, abstract and keywords at \end{frontmatter}.

\section{Introduction}
\label{sec:intro}
%%%%%%%%% INTRODUCTION %%%%%%%%%
Introduction here.
%%%%%%%%% INTRODUCTION %%%%%%%%%

\section{Computational methods}
\label{sec:methods}
%%%%%%%%% COMPUTATIONAL METHODS %%%%%%%%%
Computational methods here.
%%%%%%%%% COMPUTATIONAL METHODS %%%%%%%%%

\section{Results and discussion}
\label{sec:results}
%%%%%%%%% RESULTS AND DISCUSSION %%%%%%%%%
Results and discussion here.
%%%%%%%%% RESULTS AND DISCUSSION %%%%%%%%%

\section{Conclusions}
\label{sec:conclusions}
%%%%%%%%% CONCLUSIONS %%%%%%%%%
Conclusions here.
%%%%%%%%% CONCLUSIONS %%%%%%%%%

%%%%%%%%% BACK MATTER (unnumbered; Elsevier order) %%%%%%%%%
\section*{CRediT authorship contribution statement}
%%%%%%%%% CREDIT AUTHORSHIP CONTRIBUTION STATEMENT %%%%%%%%%
CRediT authorship contribution statement here.
%%%%%%%%% CREDIT AUTHORSHIP CONTRIBUTION STATEMENT %%%%%%%%%

\section*{Declaration of competing interest}
%%%%%%%%% DECLARATION OF COMPETING INTEREST %%%%%%%%%
Declaration of competing interest here.
%%%%%%%%% DECLARATION OF COMPETING INTEREST %%%%%%%%%

\section*{Declaration of generative AI and AI-assisted technologies in the manuscript preparation process}
%%%%%%%%% DECLARATION OF GENERATIVE AI %%%%%%%%%
Declaration of generative AI use here.
%%%%%%%%% DECLARATION OF GENERATIVE AI %%%%%%%%%

\section*{Data availability}
%%%%%%%%% DATA AVAILABILITY %%%%%%%%%
Data availability here.
%%%%%%%%% DATA AVAILABILITY %%%%%%%%%

\section*{Supplementary material}
%%%%%%%%% SUPPLEMENTARY MATERIAL %%%%%%%%%
Supplementary material here.
%%%%%%%%% SUPPLEMENTARY MATERIAL %%%%%%%%%

\section*{Acknowledgements}
%%%%%%%%% ACKNOWLEDGEMENTS %%%%%%%%%
% Elsevier: funding sources go here as well (only as stated in the materials).
Acknowledgements here.
%%%%%%%%% ACKNOWLEDGEMENTS %%%%%%%%%

%%%%%%%%% Appendix %%%%%%%%%
% If you want to add an appendix, uncomment the following lines
% (equations, figures and tables are then numbered A.1, A.2, ...):
% \appendix
% \section{xxx}
%%%%%%%%% Appendix %%%%%%%%%

%%%%%%%%% REFERENCES %%%%%%%%%
\bibliographystyle{elsarticle-num}
\bibliography{references}
%%%%%%%%% REFERENCES %%%%%%%%%

\end{document}
